

\iftrue
\section{Mechanical Systems}

This section presents a general approach for discretizing a large class of structural equilibrium problems that are governed by partial differential equations to formulate a representative discrete internal force vector $\intern$. 
The procedure generalizes the classical finite element method which is typically formulated on Sobolev vector spaces \citep{hughes1987finite} to the setting of nonlinear manifolds{\footnote{A manifold is a mathematical representation of a curved surface. Manifolds naturally arise as solution sets of constrained systems of equations and as graphs of functions.}}, which are encountered in the geometrically nonlinear analysis of plates, rods, and shells.
In all such problems, equilibrium can be stated in terms of nonlinear partial differential equations of the form:
\begin{equation}
\rho \ddot{\bm{u}} = \bm{f} -  \mathbb{B}_\chi\bm{\sigma},
\label{eq:strong}
\end{equation}
where \(\mathbb{B}_\chi\) is a nonlinear
differential operator and both the generalized stress \(\bm{\sigma}\) and the body force \(\bm{f}\) are \emph{continuous} functions defined over a body \(\mathcal{B}\) placed in space by a
map \(\chi\). In general \(\mathbb{B}\) may depend nonlinearly on \(\chi\),
but always acts linearly on \(\bm{\sigma}\). 
% This class of problems is formally summarized in Box~\ref{box:mechanical-equilibrium}.
% Problems with the form of
% \eqref{eq:strong} describe the equilibrium of continua, rods, and shells, whose
% configurations \(\chi\) may comprise a \emph{nonlinear} space
% \(\mathscr{C}\).
% Additional resources for this material are contained in
% \cite{marsden1994mathematical,zienkiewicz2014finite,ogden1997nonlinear} and \cite{shilov1996elementary}.
\iftrue

\begin{ambox}[Continuous Mechanics]{box:mechanical-equilibrium}
\begin{description}
% \tightlist
\item[Body]
A body \(\mathcal{B}\) is a set of points \(\bm{\xi}\), and a placement
\(\mathcal{B}_\chi\) of \(\mathcal{B}\) is a Riemannian manifold
\(\mathcal{B}_\chi\) with boundary \(\partial \mathcal{B}\) and tangent
bundle \(T\mathcal{B}_\chi\) embedded in an ambient Euclidean space
\(\mathcal{E}\). The map
\(\chi\colon \mathcal{B}\rightarrow \mathcal{B}_\chi\) associates the
state of a material point \(\bm{\xi}\) to the placement
\(\mathcal{B}_\chi\). A reference configuration
\(\chi_0\colon \mathcal{B}\rightarrow \mathcal{B}_{\chi_0}\) is often
implied and the pull-back by the composition \(\chi\circ\chi_0^{-1}\) is
denoted by \(\bm{A}\).
% \tightlist
\item[Energy]
The set \(\mathscr{C}\) of configurations \(\chi\) forms a smooth
 manifold. A configuration \(\chi\) is identified by a couple
\((\bm{x}, \bm{\Lambda})\) where \(\bm{x}\in\mathscr{C}_x\) is a
vectorial configuration variable and
\(\bm{\Lambda}\in\mathscr{C}_{\scriptscriptstyle{\Lambda}}\) a
non-vectorial variable. These enter independently in a Riemannian metric
on the tangent space \(T_\chi\mathscr{C}\) to \(\mathscr{C}\) at
\(\chi\), which takes the assumed form: 
\begin{equation}
  \left< \cdot , \, \cdot\right>_{\chi} = \left< \cdot , \, \cdot\right>_{x} + \left< \cdot , \, \cdot\right>_{\scriptscriptstyle{\Lambda}} .
  \label{eq:conf-pair}
\end{equation}
%
\item[Stress]
The spatial stress \(\bm{\sigma}\) in \eqref{eq:strong} is related to a
material stress \(\bm{s}\) which is defined with respect to inner products 
in \(\mathscr{C}\) at the current \(\chi\) and reference
\(\chi_0\) configurations:
\begin{equation}
  \left<\bm{c},\, \bm{\sigma}\right>_{\chi}
  =\left< \bm{A}\bm{c},\, \bm{s}\right>_{\chi_0}
  = \left<\bm{c},\, \bm{A}_*\bm{s}\right>_{\chi}
  %=\left< \bm{A}\bm{c},\, \bm{s}\right>_{\chi_0}\qquad\text{ for any }\bm{c},
  \label{eq:def-Astar}
\end{equation}
for any \(\bm{c}\).
% with \(\bm{A}_*\) denoting the Piola transformation. 
Similarly, a dual
\(\mathbb{B}^*\) to the equilibrium operator \(\mathbb{B}\) in
\eqref{eq:strong} is defined with respect to \eqref{eq:conf-pair}:
\begin{equation}
  \left<\bm{\eta},\, \mathbb{B}_\chi\bm{\sigma}\right>_\chi = \left<\mathbb{B}_\chi^*\bm{\eta},\, \bm{\sigma}\right>_\chi.
  \label{eq:def-adjoint}
\end{equation}
% \tightlist
\item[Strain]
The material strain \(\bm{e}\) is compatible with
\(\mathbb{B}^*_\chi\) through the differential relation: 
\begin{equation}
  D\bm{e}\cdot\bm{u} \equiv \bm{A}\mathbb{B}^*_\chi \bm{u}.
  \label{eq:vstrain}
\end{equation}
%
\item[Constitution]
The constitutive response is defined between the material stress
\(\bm{s}\) and the material strain \(\bm{e}\) with material \(\bm{C}_s\)  moduli defined by:
\begin{equation}
  D \bm{s}\cdot \bm{u} = \bm{C}_s \, D \bm{e}\cdot {\bm{u}}.
  \label{eq:def-moduli}
\end{equation}
\end{description}
\end{ambox}
\fi
\fi 
